%% sections/02-related-work.tex

\section{Related Work}
\label{sec:related}

\subsection{Player Typologies}

The dominant player typologies in game studies describe player motivations and
personality profiles rather than design-enabled behavioral patterns. Bartle's
four-type taxonomy~\cite{bartle1996hearts}, derived from observation of Multi-User
Dungeon (MUD) players, classifies players as Achievers (accumulating points and status),
Explorers (discovering hidden content), Socializers (interacting with other players),
and Killers (imposing on others). The taxonomy has been widely replicated, refined,
and contested; its chief limitation is that the categories describe player dispositions
that pre-exist the game rather than behaviors the game design enables.

Yee's empirical factor analysis of MMORPG motivations~\cite{yee2006motivations}
produced a three-component, ten-sub-motivation structure (Achievement: advancement,
mechanics, competition; Social: socializing, relationship, teamwork; Immersion:
discovery, role-playing, customization, escapism). This framework is more granular
and better supported empirically than Bartle's. Like Bartle's, however, it describes
why players engage, not what specific design choices enable specific behaviors.

BrainHex~\cite{nacke2011brainhex} introduced neuroscience-informed player types
(Seeker, Survivor, Daredevil, Mastermind, Conqueror, Socializer, Achiever), grounding
typology in neurobiological response patterns. The approach is methodologically
distinct from motivation surveys but converges on a similar design implication:
types describe audience segments, not design prescriptions.

The closest antecedent to the player-style concept in the typology literature is
Isbister's work on social game design~\cite{isbister2006better}, which identifies
interaction patterns enabled by game structure rather than player disposition. The
Marathon player styles extend this direction by providing a formal promotion pipeline
and cross-campaign validation evidence.

\subsection{Emergent Narrative and Emergent Behavior}

Emergent behavior in games --- player behaviors not explicitly designed --- has been
studied primarily at the level of game systems. Juul's~\cite{juul2005half} analysis
of game rules and their emergent properties establishes the foundational framework:
simple rules generate complex outcomes. Salen and Zimmerman~\cite{salen2004rules}
define emergence as the complexity arising from rule interaction and identify it as
central to meaningful play.

The Marathon player styles are not system-emergent in the Juul/Salen-Zimmerman sense.
They are design-enabled: specific character-sheet features produce specific behavioral
tendencies, which are observable as patterns across sessions. The distinction matters
because system-emergent behaviors are often difficult to predict or control by design,
whereas design-enabled behaviors are, in principle, replicable through deliberate
character-sheet choices.

Ryan's~\cite{ryan2006narrative} analysis of narrative emergence in interactive systems
identifies how stories arise from rule-governed agent interactions, a framework that
applies productively to the session-level narrative dynamics in which the Marathon
styles appear. The player styles can be understood as character-level narrative affordances:
design features that make certain narrative events (naming a grief, completing an arc)
available to the system.

\subsection{AI Player Behavior in Games}

Research on AI behavior in TRPGs and narrative games has grown alongside the field of
interactive narrative. Cavazza and colleagues~\cite{cavazza2002character} demonstrated
that AI characters in interactive drama can exhibit goal-directed narrative behavior.
Mateas and Stern's~\cite{mateas2003architecture} Fa\c{c}ade system showed that AI agents
playing within a drama manager could produce session-level narrative coherence.

More recent work on large-language-model-based game agents~\cite{akoury2023rik} has
shown that LLMs can generate contextually appropriate TRPG-style narration, though the
question of whether LLM-based agents exhibit reproducible behavioral \textit{styles}
tied to character design remains underexplored. The Marathon corpus addresses this
question directly: by constraining AI play to explicit Playstyle Heuristics on character
sheets, the corpus makes behavioral style a design input rather than an emergent
property of model behavior.

\subsection{Design Patterns in Games}

Design patterns as a methodology for game design were introduced by Bjork and
Holopainen~\cite{bjork2005patterns}, who catalogued 200+ game design patterns as
reusable problem-solution pairs. The Marathon player styles are a species of design
pattern, but at the character-sheet level rather than the game-system level: they
describe recurring behavioral solutions enabled by recurring character-sheet design
features.

Dormans~\cite{dormans2012engineering} extended the design-pattern approach to game
mechanics, and Treanor and colleagues~\cite{treanor2012game} applied it to game
narratives. Neither addresses the TRPG character-sheet as a locus of behavioral
pattern design. The present paper contributes this level of analysis.
